% ══════════════════════════════════════════════════════════════════
% CHAPITRE 2 — ANALYSE, CONCEPTION ET MODÉLISATION UML
% Plan : Intro · 2.1 · 2.2 · 2.3 · 2.4 · 2.5 · Conclusion
% ══════════════════════════════════════════════════════════════════

\chapter{Analyse, Conception et Modélisation UML}
\label{chap:conception}

\section*{Introduction}
\addcontentsline{toc}{section}{Introduction}

Après avoir positionné notre travail par rapport à l'existant, ce
chapitre aborde la phase d'analyse et de conception du système
Agro-IA Tchad. Il identifie d'abord les acteurs et les problèmes à
résoudre, spécifie ensuite les besoins fonctionnels et non
fonctionnels, présente la modélisation UML du système, décrit
l'architecture du cadre décisionnel et détaille enfin le modèle de
la base de données assurant l'historique des analyses.

% ══════════════════════════════════════════════════════════════════
\section{Analyse du problème et identification des acteurs}
\label{sec:acteurs}
% ══════════════════════════════════════════════════════════════════

\subsection{Analyse du problème}

L'analyse du contexte agricole tchadien, menée au chapitre
précédent, fait ressortir quatre problèmes opérationnels auxquels
le système doit répondre :

\begin{enumerate}[leftmargin=1.2cm]
	\item \textbf{Gestion approximative de l'irrigation} : en
	l'absence de données objectives (humidité du sol, ETP,
	prévisions), l'agriculteur irrigue selon son intuition,
	avec pour conséquences un gaspillage d'eau ou un stress
	hydrique des cultures ;
	\item \textbf{Fertilisation inadaptée} : le choix de l'engrais
	se fait rarement sur la base d'une analyse du sol,
	entraînant des dépenses inutiles et un appauvrissement
	progressif des terres ;
	\item \textbf{Détection tardive des maladies} : faute de
	diagnostic précoce, les maladies se propagent et causent
	des pertes de récolte importantes ;
	\item \textbf{Absence de traçabilité} : aucune trace des
	décisions passées n'est conservée, ce qui empêche tout
	suivi de l'évolution d'une parcelle dans le temps.
\end{enumerate}

\subsection{Identification des acteurs}

Deux acteurs interagissent avec le système :

\begin{itemize}[leftmargin=1.2cm]
	\item \textbf{L'Agriculteur (ou agent de terrain)} --- acteur
	principal : il saisit les paramètres de sa parcelle,
	consulte la météo, lance les analyses et consulte
	l'historique ;
	\item \textbf{Les services externes} --- acteurs secondaires :
	l'API Open-Meteo qui fournit les données météorologiques
	en temps réel, et la plateforme Sentinel-2/EO Browser qui
	fournit les valeurs NDVI.
\end{itemize}

% ══════════════════════════════════════════════════════════════════
\section{Spécification des besoins}
\label{sec:besoins}
% ══════════════════════════════════════════════════════════════════

\subsection{Besoins fonctionnels}

Le tableau~\ref{tab:besoins_fonc} récapitule les besoins
fonctionnels du système, regroupés par module.

\begin{table}[H]
	\centering
	\caption{Besoins fonctionnels du système Agro-IA Tchad}
	\label{tab:besoins_fonc}
	\small
	\begin{tabularx}{\textwidth}{clX}
		\toprule
		\textbf{Code} & \textbf{Module} & \textbf{Besoin fonctionnel} \\
		\midrule
		BF1 & Météo         & Afficher la météo temps réel des 5 villes
		(T°, humidité, pluie, vent, ETP) \\
		BF2 & Irrigation    & Prédire le besoin en irrigation (OUI/NON)
		avec dose estimée (mm) et confiance \\
		BF3 & Fertilisation & Recommander un engrais (Urée, DAP, NPK
		15-15-15, NPK 23-10-5, KCl) avec dose et
		période d'application \\
		BF4 & Maladies      & Diagnostiquer l'état d'une culture par
		photographie de feuille (CNN) \\
		BF5 & Maladies      & Évaluer l'état d'une parcelle par indice
		NDVI (satellite ou caméra smartphone) \\
		BF6 & Historique    & Sauvegarder automatiquement chaque analyse
		dans la base de données \\
		BF7 & Historique    & Afficher l'historique, les statistiques et
		permettre l'export CSV \\
		BF8 & Historique    & Permettre la suppression de l'historique
		par module ou en totalité \\
		\bottomrule
	\end{tabularx}
\end{table}

\subsection{Besoins non fonctionnels}

\begin{itemize}[leftmargin=1.2cm]
	\item \textbf{Performance} : temps de réponse d'une analyse
	inférieur à 3 secondes ;
	\item \textbf{Disponibilité} : accès 24h/24 depuis tout
	navigateur (ordinateur, tablette, smartphone) ;
	\item \textbf{Utilisabilité} : interface en français, simple,
	avec pré-remplissage automatique des données météo ;
	\item \textbf{Fiabilité} : accuracy minimale de 75\,\% pour les
	modules Irrigation et Fertilisation ;
	\item \textbf{Légèreté} : modèle de détection embarqué au format
	TFLite (moins de 10 Mo) pour un chargement rapide ;
	\item \textbf{Coût} : utilisation exclusive de services et
	d'outils gratuits (open source, APIs libres).
\end{itemize}

\newpage
% ══════════════════════════════════════════════════════════════════
\section{Modélisation UML du système}
\label{sec:uml}
% ══════════════════════════════════════════════════════════════════

\subsection{Diagramme de contexte}

La figure~\ref{fig:contexte} situe le système dans son
environnement : l'agriculteur est l'unique acteur humain ; les
services externes (Open-Meteo, Sentinel-2) alimentent le système
en données, et la base SQLite assure la persistance locale.

\begin{figure}[H]
	\centering
	\includegraphics[width=0.7\linewidth]{/home/tourabi/Documents/PFE_Agriculture_IA/docs/diagrammes/1_diagramme_contexte.png}
	\caption{Diagramme de contexte du système Agro-IA Tchad}
	\label{fig:contexte}
\end{figure}

\newpage
\subsection{Diagramme de cas d'utilisation}

La figure~\ref{fig:usecase} présente les cinq cas d'utilisation
offerts à l'agriculteur. Les relations
\guillemotleft~include~\guillemotright{} matérialisent le recours
automatique aux services externes : les modules d'analyse
sollicitent l'API Open-Meteo pour pré-remplir les formulaires, et
la détection par indice de végétation s'appuie sur les valeurs
NDVI relevées sur Sentinel-2.

\begin{figure}[H]
	\centering
	\includegraphics[width=0.85\linewidth]{/home/tourabi/Documents/PFE_Agriculture_IA/docs/diagrammes/cas d'utilisation.drawio.png}
	\caption{Diagramme de cas d'utilisation global avec services externes}
	\label{fig:usecase}
\end{figure}

\newpage
\subsection{Diagramme de séquence --- analyse d'une parcelle}

La figure~\ref{fig:sequence} détaille le scénario nominal d'une
analyse : pré-remplissage météo, inférence du modèle, puis
sauvegarde automatique dans l'historique avant l'affichage du
résultat.

\begin{figure}[H]
	\centering
	\includegraphics[width=0.9\linewidth]{/home/tourabi/Documents/PFE_Agriculture_IA/docs/diagrammes/3_sequence_analyse_parcelle.png}
	\caption{Diagramme de séquence --- analyse d'une parcelle}
	\label{fig:sequence}
\end{figure}

\subsection{Diagramme d'activités --- génération d'une recommandation}

La figure~\ref{fig:activites} formalise le flot d'activités commun
aux trois modules d'analyse, de la saisie des données jusqu'à
l'affichage et l'archivage de la recommandation.

\begin{figure}[H]
	\centering
	\includegraphics[width=0.65\linewidth]{/home/tourabi/Documents/PFE_Agriculture_IA/docs/diagrammes/activites.drawio.png}
	\caption{Diagramme d'activités --- génération d'une recommandation}
	\label{fig:activites}
\end{figure}

% ══════════════════════════════════════════════════════════════════
\section{Architecture du cadre décisionnel}
\label{sec:architecture}
% ══════════════════════════════════════════════════════════════════

L'architecture du système est présentée sous deux angles
complémentaires : la vue fonctionnelle en pipeline
(figure~\ref{fig:archi_generale}), qui suit le cheminement des
données depuis leurs sources jusqu'à l'agriculteur, et la vue en
couches logicielles (figure~\ref{fig:architecture}).

\begin{figure}[H]
	\centering
	\includegraphics[width=1\linewidth ]{/home/tourabi/Documents/PFE_Agriculture_IA/docs/diagrammes/architectures generales.drawio.png}
	\caption{Architecture générale du système Agro-IA Tchad}
	\label{fig:archi_generale}
\end{figure}

La vue en couches logicielles complète cette lecture :


\begin{figure}[H]
	\centering
	\includegraphics[width=0.9\linewidth]{../docs/diagrammes/architecture_couches_agro_ia}
	\caption{Architecture en couches du cadre décisionnel Agro-IA Tchad}
	\label{fig:architecture}
\end{figure}

\begin{itemize}[leftmargin=1.2cm]
	\item la \textbf{couche présentation} regroupe les cinq pages
	Streamlit (accueil, irrigation, fertilisation, maladies,
	historique) ;
	\item la \textbf{couche traitement} héberge les modèles d'IA et
	la logique métier (calcul des doses, seuils NDVI par
	culture, conseils de traitement) ;
	\item la \textbf{couche données} assure la persistance :
	modèles entraînés sérialisés et base SQLite de
	l'historique ;
	\item la \textbf{couche services externes} fournit les données
	météorologiques temps réel et les indices de végétation.
\end{itemize}

% ══════════════════════════════════════════════════════════════════
\section{Modèle de la base de données de l'historique}
\label{sec:bdd_historique}
% ══════════════════════════════════════════════════════════════════

Le module Historique répond au besoin de traçabilité (BF6 à BF8).
Le SGBD retenu est \textbf{SQLite} \cite{sqlite2024} : embarqué,
sans serveur, sans configuration, il est parfaitement adapté à
une application mono-utilisateur déployée sur le cloud. La base
\texttt{historique.db} comporte trois tables indépendantes, une
par module d'analyse.

\begin{table}[H]
	\centering
	\caption{Structure de la table \texttt{irrigation}}
	\label{tab:table_irr}
	\small
	\begin{tabularx}{\textwidth}{llX}
		\toprule
		\textbf{Attribut} & \textbf{Type} & \textbf{Description} \\
		\midrule
		id             & INTEGER & Clé primaire auto-incrémentée \\
		date\_analyse  & TEXT    & Date et heure de l'analyse \\
		ville, culture, stade & TEXT & Localisation et culture analysée \\
		zone, saison, type\_sol & TEXT & Contexte agro-climatique \\
		temperature\_moy, humidite\_air & REAL & Conditions climatiques \\
		pluie\_7j, humidite\_sol, etp   & REAL & Variables hydriques \\
		resultat       & TEXT    & Décision : OUI / NON \\
		dose\_mm       & REAL    & Dose d'irrigation estimée (mm) \\
		confiance      & REAL    & Probabilité associée (\%) \\
		\bottomrule
	\end{tabularx}
\end{table}

\begin{table}[H]
	\centering
	\caption{Structure de la table \texttt{fertilisation}}
	\label{tab:table_fert}
	\small
	\begin{tabularx}{\textwidth}{llX}
		\toprule
		\textbf{Attribut} & \textbf{Type} & \textbf{Description} \\
		\midrule
		id             & INTEGER & Clé primaire auto-incrémentée \\
		date\_analyse  & TEXT    & Date et heure de l'analyse \\
		ville, culture, type\_sol, zone & TEXT & Contexte de la parcelle \\
		azote\_n, phosphore\_p, potassium\_k & REAL & Analyse NPK du sol \\
		ph\_sol        & REAL    & Acidité du sol \\
		temperature, humidite\_air, pluie & REAL & Conditions climatiques \\
		engrais\_recommande & TEXT & Engrais prédit par XGBoost \\
		confiance      & REAL    & Probabilité associée (\%) \\
		\bottomrule
	\end{tabularx}
\end{table}

\begin{table}[H]
	\centering
	\caption{Structure de la table \texttt{maladies}}
	\label{tab:table_mal}
	\small
	\begin{tabularx}{\textwidth}{llX}
		\toprule
		\textbf{Attribut} & \textbf{Type} & \textbf{Description} \\
		\midrule
		id             & INTEGER & Clé primaire auto-incrémentée \\
		date\_analyse  & TEXT    & Date et heure de l'analyse \\
		ville, culture & TEXT    & Contexte de la parcelle \\
		methode        & TEXT    & CNN Photo / NDVI satellite /
		NDVI caméra \\
		resultat, etat & TEXT    & Classe prédite et état
		(Saine, Stress, Malade) \\
		confiance      & REAL    & Probabilité associée (\%) \\
		ndvi           & REAL    & Valeur NDVI (le cas échéant) \\
		temperature, humidite\_air, pluie\_7j & REAL & Contexte
		climatique \\
		traitement     & TEXT    & Conseil de traitement fourni \\
		\bottomrule
	\end{tabularx}
\end{table}

Ce modèle volontairement dénormalisé (chaque table est autonome,
sans clé étrangère) simplifie les insertions et les lectures,
supprime tout risque d'incohérence référentielle et facilite
l'export CSV table par table --- des choix cohérents avec un
usage mono-utilisateur et un volume de données modéré.

\section*{Conclusion}
\addcontentsline{toc}{section}{Conclusion}

Ce chapitre a formalisé la conception du système Agro-IA Tchad :
identification des acteurs et des problèmes, spécification de
huit besoins fonctionnels couvrant les quatre modules,
modélisation UML des interactions, architecture en quatre couches
et modèle de données de l'historique. Cette conception constitue
le socle de l'implémentation présentée au chapitre suivant.